\section{Overview and Problem Formulation}
\label{section:overview}

\subsection{Motivation: Addressing the Parameter-Efficiency Bottleneck in ResMLP}
At the conclusion of our Week 02 benchmark, our team identified a fundamental structural limitation in columnar vision MLP architectures. Under identical experimental conditions on CIFAR-100, vanilla ResMLP~\cite{touvron2021resmlp} required 2,191,204 parameters to achieve a Top-1 validation accuracy of 62.40\%. In contrast, modern hierarchical architectures such as PoolFormer~\cite{yu2022metaformer} and RepMLP~\cite{ding2022repmlp} demonstrated substantially higher accuracy with equal or fewer parameters.

This performance gap raised an important research question: \textit{Why does columnar ResMLP struggle to learn efficient visual representations from small datasets, and why do hierarchical architectures achieve significantly higher parameter efficiency?}

Our team formulated the working hypothesis that isotropic columnar models suffer from two structural disadvantages:
\begin{itemize}
    \item \textbf{Absence of multi-scale feature hierarchies:} In vanilla ResMLP, non-overlapping patch slicing maintains a rigid, single-scale token sequence ($N = 64$ for patch size $P = 4$) throughout all network layers. Because spatial token dimensions never contract, early layers are forced to process global token interactions prematurely, while later layers lack compact, semantically rich spatial representations.
    \item \textbf{Lack of 2D translation equivariance:} A dense linear token mixing matrix $\mathbf{W} \in \mathbb{R}^{N \times N}$ treats all spatial positions as arbitrary permutation-invariant tokens. It does not possess any built-in concept of adjacent 2D pixel proximity, requiring massive training data or heavy regularizations to learn spatial relationships from scratch.
\end{itemize}

We hypothesized that introducing convolutional patch downsampling (a pyramid backbone) before each stage shrinks the token sequence, broadens the effective receptive field, and provides strong multi-scale inductive bias, enabling the network to achieve high accuracy with a fraction of the parameter cost.

\subsection{Research Objectives for Week 03}
To verify this hypothesis and explore the broader design space of spatial mixing, our team conducted five complementary research investigations in Week 03:
\begin{itemize}
    \item \textbf{a. Validating hierarchical pyramids with re-parameterization:} We developed Pyramid-ResMLP, which integrates an overlapping convolution stem, a 3-stage feature pyramid ($16 \times 16 \rightarrow 8 \times 8 \rightarrow 4 \times 4$ token resolution for the $32 \times 32$ CIFAR-100 dataset), and a structural re-parameterization token mixer (RepTokenMix layer). To isolate the exact source of performance gains, we also trained Pyramid-ResMLP Pure as an unaugmented ablation control.
    \item \textbf{b. Evaluating parameter-free spatial communication:} We evaluated Shift-ResMLP, which replaces matrix multiplication in token mixing with a 4-direction coordinate channel shift. This eliminates all learnable spatial weights, providing an empirical test of whether non-parametric coordinate routing can provide sufficient inductive bias.
    \item \textbf{c. Resolution-agnostic cellular automata mixing:} We explored CA-Mixer, which replaces fixed-dimension matrix mixing with Neural Cellular Automata (NCA) updates. Because its perception operators are local 2D convolutions, its parameter count remains constant across arbitrary input resolutions.
    \item \textbf{d. MetaFormer and hybrid paradigms:} We examined non-parametric average pooling (PoolFormer-S12) and a hybrid architecture (HybridFormer) combining early-stage pooling with late-stage self-attention on CIFAR-100.
    \item \textbf{e. Scaling to Tiny-ImageNet ($64 \times 64$):} We evaluated whether the multi-stage pyramid design scales effectively from CIFAR-100 to Tiny-ImageNet (200 classes) under different channel capacities and regularization schedules.
\end{itemize}
